\section{Related work}
\label{sec:related}

\paragraph{Choosing items from source-model results.}
Anchor Points~\citep{vivek2024anchor} cluster the response patterns of existing models and keep one representative item per cluster, estimating accuracy from these ``anchors''; the authors report that a handful of anchors can rank models better than a uniform random subset of the same size. tinyBenchmarks~\citep{mauricio2024tiny} combines such anchor items with item response theory (IRT) to estimate a new model's score, and reports that about a hundred curated MMLU items can suffice. metabench~\citep{kipnis2025metabench} selects items by IRT information fitted on a large population of models and reports reconstructing benchmark scores from under 3\% of the items. The learned-assessor approach of \citet{pacchiardi2024100} predicts a new model's success on unseen instances from its results on a few reference instances, and finds that random reference instances are close to more elaborate selection rules; their assessors use instance features, which are not available here. Active Evaluation Acquisition~\citep{li2025active} learns a policy that picks evaluation items adaptively, using a model of the dependencies between items to predict the remainder. We have read these works at the level of their abstracts and do not restate their protocols; the question our protocol is built around is whether the advantage over random sampling survives when the new model is stronger than all source models.

\paragraph{Benchmark-prediction evaluations.}
\citet{zhang2025benchmark} compare eleven benchmark-prediction methods on a collection of benchmarks and release the item-level score matrices and method implementations that we use as data (Section~\ref{sec:setup}). We use their data and, where stated, read their code, and we run every method ourselves under the protocol below; no number of theirs is reused.

\paragraph{Statistical background.}
Estimating a finite-population mean from a sample with the help of a model fitted on auxiliary information is the subject of model-assisted survey sampling~\citep{sarndal1992model}: stratified sampling uses auxiliary information in the design, and the difference (residual-correction) estimator adds to the model's population mean the sample mean of its residuals, which makes the estimator design-unbiased whatever the model. The same construction underlies prediction-powered inference~\citep{angelopoulos2023ppi}. Our P3 estimator (Section~\ref{sec:methods}) is this standard estimator with a low-rank model of the sources; we do not claim it as new, only test it as a design for benchmark prediction. Item response theory~\citep{embretson2000irt} is the latent-trait model on which the IRT pipelines are based. Adaptive evaluation with continuous scores~\citep{balkir2026confident}, which does not apply to binary item results, is outside our comparison. We also do not use sample-level benchmarking over open-ended capabilities~\citep{ghosh2024onebench}, since it is not a binary item matrix of the kind supplied; MMLU-Pro~\citep{wang2024mmlupro} enters only as one of the benchmarks of the Open LLM Leaderboard v2.

\paragraph{Benchmarks and score matrices.}
The Open LLM Leaderboard v2~\citep{fourrier2024openllm} and HELM~\citep{liang2023helm} provide the item-level correctness of many models on a common set of benchmarks, which is what makes source-model-based selection possible.
